\documentclass[UTF8,12pt]{ctexart}
\usepackage[a4paper,margin=2.3cm]{geometry}
\usepackage{amsmath,amssymb}
\usepackage{array}
\usepackage{enumitem}
\usepackage{fancyhdr}
\usepackage{lastpage}
\usepackage[colorlinks,
linkcolor=black,
anchorcolor=black,
citecolor=black
]{hyperref}

\setlength{\parindent}{2em}
\setlength{\parskip}{0.35em}
\linespread{1.15}
\pagestyle{fancy}
\fancyhf{}
\cfoot{第 \thepage\ 页，共 \pageref{LastPage} 页}
\renewcommand{\headrulewidth}{0pt}

\newcommand{\blank}[1][4em]{\underline{\hspace{#1}}}
\newcommand{\examsection}[1]{\par\addvspace{1em}\noindent{\large\bfseries #1}\par\addvspace{0.4em}}
\newcommand{\choice}[4]{%
  \begin{tabular}{@{}p{0.48\linewidth}p{0.48\linewidth}@{}}
  (A) #1 & (B) #2\\[0.4em]
  (C) #3 & (D) #4
  \end{tabular}
}

\begin{document}

\begin{center}
  {\LARGE\bfseries 《概率论与数理统计》期中考试试题}\\[0.5em]
  {\large 2025--2026 学年 第二学期 / 福建师范大学数学与统计学院}\\[0.3em]
  排版：\href{https://github.com/Xuuyuan}{@许愿} -- \href{https://wefjnu.nekoark.com}{WeFJNU}
\end{center}

\examsection{一、单项选择题（每小题 3 分，共 15 分）}
\begin{enumerate}[label={\arabic*、},leftmargin=2em,itemsep=0.8em]
\item 设 $A,B$ 是两个事件，则 ``$P(A)+P(B)=1$'' 是 ``$A$ 与 $B$ 互为对立事件'' 的（\quad）。

\choice{充分非必要条件；}{必要非充分条件；}{充要条件；}{既不是充分条件，也不是必要条件。}

\item 甲、乙、丙三人各射击一次，$A,B,C$ 分别表示甲、乙、丙击中目标，则事件 ``三人中至少有一人未击中'' 可表示为（\quad）。

\begin{tabular}{@{}p{0.48\linewidth}p{0.48\linewidth}@{}}
(A) $\overline{ABC}$； &
(B) $\overline A\,\overline B\,\overline C$；\\[0.4em]
(C) $\overline A B C\cup A\overline B C\cup AB\overline C$； &
(D) $\overline{A\cup B\cup C}$。
\end{tabular}

\item 设 $A$ 和 $B$ 为随机事件，$P(AB)=0$，则下列命题中正确的是（\quad）。

\choice{$A$ 和 $B$ 互不相容；}{$AB$ 是不可能事件；}{$AB$ 未必是不可能事件；}{$P(A)=0$ 或 $P(B)=0$。}

\item 设样本空间 $S=\{1,2,3,4\}$，且每个样本点出现的概率相等，令
$A_1=\{1,2\}$，$A_2=\{1,3\}$，$A_3=\{1,4\}$，$A_4=\{2,3\}$，则下列结论正确的是（\quad）。

\begin{enumerate}[label=(\Alph*),leftmargin=2em,itemsep=0.2em]
\item $A_1,A_2,A_3$ 相互独立，$A_1,A_2,A_4$ 相互独立；
\item $A_1,A_2,A_3$ 相互独立，$A_1,A_2,A_4$ 两两独立；
\item $A_1,A_2,A_3$ 两两独立，$A_1,A_2,A_4$ 相互独立；
\item $A_1,A_2,A_3$ 两两独立，$A_1,A_2,A_4$ 两两独立。
\end{enumerate}

\item 设随机变量 $X$ 的分布函数为
\[
F(x)=
\begin{cases}
0, & x<0,\\
0.2, & 0\le x<1,\\
0.5, & 1\le x<2,\\
1, & x\ge 2,
\end{cases}
\]
则 $P\left(X=\frac12\right)=$（\quad）。

\choice{$0$；}{$0.2$；}{$0.5$；}{$a$（$a$ 为 $[0,1]$ 中任意实数）。}
\end{enumerate}

\newpage
\examsection{二、填空题（每小题 3 分，共 15 分）}
\begin{enumerate}[label={\arabic*、},leftmargin=2em,itemsep=0.8em]
\item 设随机事件 $A,B$ 相互独立，已知 $P(A-B)=P(B-A)=0.25$，则
$P(AB\mid A\cup B)=$ \blank[8em]。

\item 若离散型随机变量 $X$ 和 $Y$ 独立同分布，且 $X$ 的分布律为
\[
\begin{array}{|c|c|c|c|}
\hline
X & -2 & 0 & 2\\ \hline
p & 0.3 & 0.5 & 0.2\\ \hline
\end{array}
\]
则 $P(X=Y)$ 的值为 \blank[8em]。

\item 某人向同一目标独立重复射击，直到连续击中目标两次停止。假设每次射击击中目标的概率为
$p(0<p<1)$，则该人共射击 $5$ 次的概率为 \blank[8em]。

\item 设随机变量 $X\sim N(\mu,\sigma^2)$，$P(X<a)=0.32$，$P(X\ge 2)=0.5$，则
$P(a<X<4-a)=$ \blank[8em]。

\item 已知随机变量 $X$ 服从 Poisson 分布，且 $P(X=2)=3P(X=4)$，则
$P(X\ge 1)=$ \blank[8em]。
\end{enumerate}

\examsection{三、（8 分）}
随机变量 $X$ 服从参数为 $1$ 的指数分布，对 $X$ 独立重复观察 $5$ 次，求至少有 $1$ 次观测值大于 $3$ 的概率。

\examsection{四、（12 分）}
在 2026 年央视春晚舞台上，多款智能机器人协同完成舞蹈、列队、翻转等高难度表演。某实验室为测试 $A,B$ 两种型号机器人的动作稳定性，设计如下试验：

每次独立执行一个动作，若某型号机器人试验成功，则下一轮继续使用该型号机器人进行试验；若试验失败，则下一轮更换另外一种型号的机器人进行试验。

已知 $A$ 型号机器人试验成功的概率为 $\frac23$，$B$ 型号机器人试验成功的概率为 $\frac12$。每次试验成功记 $1$ 分，失败记 $0$ 分，且第 $1$ 轮使用 $A$ 型号机器人。

记 $X$ 为前 $3$ 轮试验的总得分，求 $X$ 的分布律。

\examsection{五、（12 分）}
某同学参加学校组织的数学知识竞赛。在 $5$ 道四选一的单选题中有 $3$ 道有思路，有 $2$ 道完全没有思路，有思路的题目每道做对的概率为 $\frac12$，而没有思路的题目只好任意猜一个答案。若从这 $5$ 道题目中任选 $2$ 题，该同学 $2$ 道题目都做对了，求 $2$ 道题都是猜对的概率。

\newpage
\examsection{六、（12 分）}
设 $(X,Y)$ 的联合密度函数为
\[
f(x,y)=
\begin{cases}
6x, & 0<x<y<1,\\
0, & \text{其他}.
\end{cases}
\]
\begin{enumerate}[label=(\arabic*),leftmargin=2em,itemsep=0.6em]
\item 求边缘密度函数 $f_X(x),f_Y(y)$；
\item 求 $P(X+Y\le 1)$。
\end{enumerate}

\examsection{七、（14 分）}
设随机变量 $X$ 的概率密度为
\[
f(x)=
\begin{cases}
x, & 0\le x<1,\\
c-x, & 1\le x<2,\\
0, & \text{其他}.
\end{cases}
\]
\begin{enumerate}[label=(\arabic*),leftmargin=2em,itemsep=0.6em]
\item 求常数 $c$；
\item 求分布函数 $F(x)$；
\item 求 $P(X=1)$。
\end{enumerate}

\examsection{八、（12 分）}
设随机变量 $X\sim U(0,2)$，求随机变量 $Y=|\ln X|$ 的概率密度函数。

\end{document}
